\documentclass{article}
\usepackage{amsmath,amssymb,amsthm}
\usepackage{algorithm}
\usepackage{algpseudocode}
\usepackage{hyperref}
\usepackage{enumitem}
\newtheorem{theorem}{Theorem}
\newtheorem{lemma}{Lemma}
\newtheorem{assumption}{Assumption}
\newcommand{\E}{\mathbb{E}}
\newcommand{\Var}{\mathrm{Var}}
\newcommand{\R}{\mathbb{R}}

\begin{document}

\title{Stochastic Gradient Langevin Dynamics with Decaying Noise\\
(Non‑convex Smooth Optimization)}
\author{}
\date{}
\maketitle

%%%%%%%%%%%%%%%%%%%%%%%%%%%%%%%%%%%%%%%%%%%%%%%%%%%%%%%%%%%%%%%%%
\section*{Algorithm Name}
\textbf{Stochastic Gradient Langevin Dynamics (SGLD) with Time‑varying Noise}\\
The method combines a standard gradient step with an additive Gaussian
perturbation whose variance is scheduled to decrease over iterations.

%%%%%%%%%%%%%%%%%%%%%%%%%%%%%%%%%%%%%%%%%%%%%%%%%%%%%%%%%%%%%%%%%
\section*{Mathematical Formulation}
Let $f:\R^{d}\rightarrow\R$ be the objective function.
Given a stepsize (learning rate) $\alpha>0$ and a sequence of noise
covariances $\{\Sigma_{k}\}_{k\ge 0}$,
the iterates $\{x_{k}\}_{k\ge 0}$ are generated by  

\[
\boxed{
x_{k+1}=x_{k}-\alpha\,\nabla f(x_{k})+\xi_{k}},\qquad k=0,1,\dots
\tag{1}
\]

where  

\[
\xi_{k}\sim\mathcal N(0,\Sigma_{k}),\qquad 
\Sigma_{k}= \sigma_{k}^{2} I_{d},\qquad 
\sigma_{k}^{2}= \frac{c}{k+1},\;c>0.
\tag{2}
\]

Thus the noise is unbiased,
\[
\E[\xi_{k}\mid x_{k}]=0,\qquad \E[\|\xi_{k}\|^{2}\mid x_{k}]=d\sigma_{k}^{2},
\tag{3}
\]
and its variance decays as $O(1/k)$.

%%%%%%%%%%%%%%%%%%%%%%%%%%%%%%%%%%%%%%%%%%%%%%%%%%%%%%%%%%%%%%%%%
\section*{Pseudocode}
\begin{algorithm}[H]
\caption{SGLD with Decaying Noise}
\begin{algorithmic}[1]
\Require Initial point $x_{0}\in\R^{d}$, stepsize $\alpha>0$, constant $c>0$, maximum iterations $T$.
\For{$k=0$ \textbf{to} $T-1$}
    \State Compute (exact) gradient $g_{k}= \nabla f(x_{k})$.
    \State Set noise variance $\sigma_{k}^{2}=c/(k+1)$.
    \State Sample $\xi_{k}\sim\mathcal N(0,\sigma_{k}^{2} I_{d})$.
    \State Update 
    \[
    x_{k+1}=x_{k}-\alpha g_{k}+ \xi_{k}.
    \]
\EndFor
\State \textbf{output} $x_{T}$ (or a random iterate uniformly chosen from $\{x_{0},\dots,x_{T-1}\}$).
\end{algorithmic}
\end{algorithm}

%%%%%%%%%%%%%%%%%%%%%%%%%%%%%%%%%%%%%%%%%%%%%%%%%%%%%%%%%%%%%%%%%
\section*{Dry Run Example}
Consider the one‑dimensional quadratic  

\[
f(w)=(w-3)^{2},\qquad \nabla f(w)=2(w-3).
\]

Parameters: $w_{0}=0$, $\alpha=0.1$, $c=1$ (so $\sigma_{k}^{2}=1/(k+1)$).  
For illustration we set the random seed so that $\xi_{0}=0.05$, $\xi_{1}=-0.03$, $\xi_{2}=0.02$.

\begin{enumerate}[label=Iteration \arabic*:, leftmargin=*]
\item $k=0$  
\[
g_{0}=2(w_{0}-3)=-6,\qquad 
w_{1}=w_{0}-\alpha g_{0}+\xi_{0}=0-0.1(-6)+0.05=0.65.
\]  
Objective: $f(w_{1})=(0.65-3)^{2}=5.5225$.

\item $k=1$  
\[
g_{1}=2(0.65-3)=-4.7,\qquad 
\sigma_{1}^{2}=1/2,\; \xi_{1}=-0.03,
\]  
\[
w_{2}=0.65-0.1(-4.7)-0.03=1.12.
\]  
Objective: $f(w_{2})=(1.12-3)^{2}=3.5344$.

\item $k=2$  
\[
g_{2}=2(1.12-3)=-3.76,\qquad 
\sigma_{2}^{2}=1/3,\; \xi_{2}=0.02,
\]  
\[
w_{3}=1.12-0.1(-3.76)+0.02=1.516.
\]  
Objective: $f(w_{3})=(1.516-3)^{2}=2.197.
\]
\end{enumerate}
The iterates move toward the minimiser $w^{*}=3$ while the injected noise
diminishes.

%%%%%%%%%%%%%%%%%%%%%%%%%%%%%%%%%%%%%%%%%%%%%%%%%%%%%%%%%%%%%%%%%
\section*{Assumptions}
\begin{assumption}[Smoothness]\label{ass:smooth}
$f$ is $L$‑smooth, i.e. for all $x,y\in\R^{d}$
\[
\|\nabla f(x)-\nabla f(y)\|\le L\|x-y\| .
\tag{A1}
\]
Equivalently,
\[
f(y)\le f(x)+\langle\nabla f(x),y-x\rangle+\frac{L}{2}\|y-x\|^{2}.
\tag{A1'}
\]
\end{assumption}

\begin{assumption}[Lower Boundedness]\label{ass:lower}
There exists $f^{*}>-\infty$ such that $f(x)\ge f^{*}$ for all $x$.
\tag{A2}
\end{assumption}

\begin{assumption}[Noise]\label{ass:noise}
For each $k\ge0$,
\[
\E[\xi_{k}\mid x_{k}]=0,\qquad 
\E[\|\xi_{k}\|^{2}\mid x_{k}]=d\sigma_{k}^{2},
\quad\sigma_{k}^{2}= \frac{c}{k+1},\;c>0 .
\tag{A3}
\]
The sequence $\{\sigma_{k}^{2}\}$ is non‑increasing and satisfies
\[
\sum_{k=0}^{\infty}\sigma_{k}^{2}<\infty .
\tag{A3'}
\]
\end{assumption}

\begin{assumption}[Stepsize]\label{ass:stepsize}
The learning rate $\alpha$ is constant and chosen such that
\[
0<\alpha\le \frac{1}{L}.
\tag{A4}
\]
\end{assumption}

%%%%%%%%%%%%%%%%%%%%%%%%%%%%%%%%%%%%%%%%%%%%%%%%%%%%%%%%%%%%%%%%%
\section*{Main Convergence Theorem}
\begin{theorem}[Average Gradient Norm]\label{thm:main}
Under Assumptions \ref{ass:smooth}–\ref{ass:stepsize},
for any integer $T\ge1$ the iterates generated by \eqref{1} satisfy  

\[
\frac{1}{T}\sum_{k=0}^{T-1}\E\!\left[\|\nabla f(x_{k})\|^{2}\right]
\;\le\;
\frac{2\bigl(f(x_{0})-f^{*}\bigr)}{\alpha T}
\;+\;
L\alpha\,\frac{d\,c}{T}\sum_{k=0}^{T-1}\frac{1}{k+1}.
\tag{4}
\]

Consequently, using the harmonic sum bound $\sum_{k=0}^{T-1}\frac{1}{k+1}
\le 1+\log T$, we obtain  

\[
\frac{1}{T}\sum_{k=0}^{T-1}\E\!\left[\|\nabla f(x_{k})\|^{2}\right]
\;\le\;
\underbrace{\frac{2\bigl(f(x_{0})-f^{*}\bigr)}{\alpha}}_{\displaystyle O(1)}
\frac{1}{T}
\;+\;
\underbrace{L\alpha d c}_{\displaystyle O(1)}\frac{1+\log T}{T}.
\tag{5}
\]

Thus the averaged squared gradient norm decays as $O\!\bigl(\frac{\log T}{T}\bigr)$.
\end{theorem}

%%%%%%%%%%%%%%%%%%%%%%%%%%%%%%%%%%%%%%%%%%%%%%%%%%%%%%%%%%%%%%%%%
\section*{Theoretical Properties}
\begin{enumerate}[label=\arabic*.]
\item \textbf{Stability.} The condition $\alpha\le1/L$ guarantees that the
deterministic part of the update is a non‑expansive map; the added
zero‑mean noise does not destabilise the recursion because its total
energy $\sum_{k}\sigma_{k}^{2}$ is finite (Assumption \ref{ass:noise} A3').

\item \textbf{Hyper‑parameter Sensitivity.}
    \begin{itemize}
        \item Larger $\alpha$ accelerates the $1/(\alpha T)$ term but
        inflates the second term $L\alpha d c$, so an optimal constant
        stepsize is $\alpha^{\star}= \Theta\!\bigl(1/\sqrt{T}\bigr)$,
        yielding $O(1/\sqrt{T})$ in the worst case.
        \item The decay constant $c$ scales the variance contribution linearly.
    \end{itemize}

\item \textbf{Comparison to Baselines.}
    \begin{itemize}
        \item \emph{Pure SGD} (no noise) obtains the same $O(1/(\alpha T))$
        bound on the gradient norm for smooth non‑convex functions.
        \item \emph{Langevin dynamics with constant variance}
        only guarantees convergence to a stationary distribution and
        typically yields a $O(1/\sqrt{T})$ bound; the decaying schedule
        restores a diminishing‑gradient guarantee.
    \end{itemize}
\end{enumerate}

%%%%%%%%%%%%%%%%%%%%%%%%%%%%%%%%%%%%%%%%%%%%%%%%%%%%%%%%%%%%%%%%%
\section*{Preliminaries (Definitions and Lemmas)}
\begin{lemma}[Smoothness Descent Inequality]\label{lem:smooth}
If $f$ is $L$‑smooth, then for any $x,y\in\R^{d}$
\[
f(y)\le f(x)+\langle\nabla f(x),y-x\rangle+\frac{L}{2}\|y-x\|^{2}.
\tag{L1}
\]
\end{lemma}
\begin{proof}
Standard consequence of $L$‑smoothness; see e.g. Nesterov (2004). $\square$
\end{proof}

\begin{lemma}[Bound on Cross Term]\label{lem:cross}
For any random vector $\xi$ with $\E[\xi]=0$ and any deterministic $a$,
\[
\E\bigl[\langle a,\xi\rangle\bigr]=0.
\tag{L2}
\]
\end{lemma}
\begin{proof}
Linearity of expectation and zero‑mean property. $\square$
\end{proof}

%%%%%%%%%%%%%%%%%%%%%%%%%%%%%%%%%%%%%%%%%%%%%%%%%%%%%%%%%%%%%%%%%
\section*{Main Proof (Step‑by‑Step Derivation)}
We prove Theorem \ref{thm:main}.  All expectations are taken w.r.t.
the randomness up to iteration $k$ unless otherwise stated.

\noindent\textbf{Step 1. Apply smoothness inequality (Lemma~\ref{lem:smooth}) with
$y=x_{k+1}$ and $x=x_{k}$.}
\[
f(x_{k+1})\le f(x_{k})+\langle\nabla f(x_{k}),x_{k+1}-x_{k}\rangle
+\frac{L}{2}\|x_{k+1}-x_{k}\|^{2}.
\tag{S1}
\]

\noindent\textbf{Step 2. Substitute the update rule \eqref{1}.}
Recall $x_{k+1}-x_{k}= -\alpha\nabla f(x_{k})+\xi_{k}$. Hence  

\[
\begin{aligned}
\langle\nabla f(x_{k}),x_{k+1}-x_{k}\rangle
&= \langle\nabla f(x_{k}),-\alpha\nabla f(x_{k})+\xi_{k}\rangle \\
&= -\alpha\|\nabla f(x_{k})\|^{2}
   +\langle\nabla f(x_{k}),\xi_{k}\rangle .
\end{aligned}
\tag{S2}
\]

\noindent\textbf{Step 3. Expand the squared norm term.}
\[
\|x_{k+1}-x_{k}\|^{2}
= \|-\alpha\nabla f(x_{k})+\xi_{k}\|^{2}
= \alpha^{2}\|\nabla f(x_{k})\|^{2}
  -2\alpha\langle\nabla f(x_{k}),\xi_{k}\rangle
  +\|\xi_{k}\|^{2}.
\tag{S3}
\]

\noindent\textbf{Step 4. Insert \eqref{S2} and \eqref{S3} into \eqref{S1}.}
\[
\begin{aligned}
f(x_{k+1}) \le &
f(x_{k})
-\alpha\|\nabla f(x_{k})\|^{2}
+\langle\nabla f(x_{k}),\xi_{k}\rangle \\
&\quad+\frac{L}{2}\Bigl(\alpha^{2}\|\nabla f(x_{k})\|^{2}
-2\alpha\langle\nabla f(x_{k}),\xi_{k}\rangle
+\|\xi_{k}\|^{2}\Bigr).
\end{aligned}
\tag{S4}
\]

\noindent\textbf{Step 5. Collect like terms.}
\[
\begin{aligned}
f(x_{k+1}) \le &
f(x_{k})
-\Bigl(\alpha-\frac{L\alpha^{2}}{2}\Bigr)\|\nabla f(x_{k})\|^{2} \\
&\;+\Bigl(1-\!L\alpha\Bigr)\langle\nabla f(x_{k}),\xi_{k}\rangle
+\frac{L}{2}\|\xi_{k}\|^{2}.
\end{aligned}
\tag{S5}
\]

\noindent\textbf{Step 6. Take conditional expectation $\E[\cdot\mid x_{k}]$.}
Using Lemma~\ref{lem:cross} (Step \ref{lem:cross}) and Assumption \ref{ass:noise}
(A3),

\[
\begin{aligned}
\E\!\bigl[f(x_{k+1})\mid x_{k}\bigr]
\le &
f(x_{k})
-\Bigl(\alpha-\frac{L\alpha^{2}}{2}\Bigr)\|\nabla f(x_{k})\|^{2}
+\frac{L}{2}\E\!\bigl[\|\xi_{k}\|^{2}\mid x_{k}\bigr] \\
=&
f(x_{k})
-\Bigl(\alpha-\frac{L\alpha^{2}}{2}\Bigr)\|\nabla f(x_{k})\|^{2}
+\frac{L}{2}d\sigma_{k}^{2}.
\end{aligned}
\tag{S6}
\]

\noindent\textbf{Step 7. Use the stepsize bound $\alpha\le 1/L$ to simplify
the coefficient.}
Since $\alpha\le1/L$, we have $L\alpha^{2}\le\alpha$, thus  

\[
\alpha-\frac{L\alpha^{2}}{2}\;\ge\;\frac{\alpha}{2}.
\tag{S7}
\]

\noindent\textbf{Step 8. Apply (S7) to (S6).}
\[
\E\!\bigl[f(x_{k+1})\mid x_{k}\bigr]
\le
f(x_{k})-\frac{\alpha}{2}\|\nabla f(x_{k})\|^{2}
+\frac{L d}{2}\sigma_{k}^{2}.
\tag{S8}
\]

\noindent\textbf{Step 9. Take full expectation and sum over $k=0,\dots,T-1$.}
Denote $\Delta_{k}:=\E[f(x_{k})]$. Then (S8) yields  

\[
\Delta_{k+1}\le\Delta_{k}-\frac{\alpha}{2}\E\!\bigl[\|\nabla f(x_{k})\|^{2}\bigr]
+\frac{L d}{2}\sigma_{k}^{2}.
\tag{S9}
\]

Summing telescopically,

\[
\Delta_{T}\le\Delta_{0}
-\frac{\alpha}{2}\sum_{k=0}^{T-1}\E\!\bigl[\|\nabla f(x_{k})\|^{2}\bigr]
+\frac{L d}{2}\sum_{k=0}^{T-1}\sigma_{k}^{2}.
\tag{S10}
\]

\noindent\textbf{Step 10. Rearrange to isolate the gradient sum.}
Because $f$ is lower bounded by $f^{*}$ (Assumption \ref{ass:lower}),
$\Delta_{T}\ge f^{*}$. Hence  

\[
\frac{\alpha}{2}\sum_{k=0}^{T-1}\E\!\bigl[\|\nabla f(x_{k})\|^{2}\bigr]
\le \Delta_{0}-f^{*}
+\frac{L d}{2}\sum_{k=0}^{T-1}\sigma_{k}^{2}.
\tag{S11}
\]

Dividing by $\alpha T$ gives the average bound (4) of Theorem~\ref{thm:main}.

\noindent\textbf{Step 11. Plug the explicit variance schedule.}
With $\sigma_{k}^{2}=c/(k+1)$,
\[
\sum_{k=0}^{T-1}\sigma_{k}^{2}=c\sum_{k=0}^{T-1}\frac{1}{k+1}
\le c\bigl(1+\log T\bigr).
\tag{S12}
\]

Insert (S12) into (S11) and simplify to obtain (5).  $\square$

%%%%%%%%%%%%%%%%%%%%%%%%%%%%%%%%%%%%%%%%%%%%%%%%%%%%%%%%%%%%%%%%%
\section*{Rate Derivation (With Constants)}
From (4) we have  

\[
\frac{1}{T}\sum_{k=0}^{T-1}\E\!\bigl[\|\nabla f(x_{k})\|^{2}\bigr]
\le
\underbrace{\frac{2\bigl(f(x_{0})-f^{*}\bigr)}{\alpha}}_{\displaystyle C_{1}}
\frac{1}{T}
\;+\;
\underbrace{L\alpha d c}_{\displaystyle C_{2}}\frac{1+\log T}{T}.
\]

Define $C_{1}=2\bigl(f(x_{0})-f^{*}\bigr)/\alpha$ and $C_{2}=L\alpha d c$.
Thus the rate is  

\[
\boxed{
\frac{1}{T}\sum_{k=0}^{T-1}\E\!\bigl[\|\nabla f(x_{k})\|^{2}\bigr]
\le \frac{C_{1}+C_{2}(1+\log T)}{T}
}
\tag{R1}
\]

which is $O(\log T/T)$.

If one selects a diminishing stepsize $\alpha_{k}=a/(k+1)^{\beta}$ with
$\beta\in(0,1]$, a similar calculation (omitted for brevity) yields the
classical $O(1/\sqrt{T})$ bound for $\beta=1/2$ and $O(1/T)$ for $\beta=1$.

%%%%%%%%%%%%%%%%%%%%%%%%%%%%%%%%%%%%%%%%%%%%%%%%%%%%%%%%%%%%%%%%%
\section*{Special Cases}
\begin{enumerate}[label=(\alph*)]
\item \textbf{Zero Noise ($c=0$).} The bound reduces to the standard
deterministic gradient descent guarantee
$\frac{1}{T}\sum_{k}\E\|\nabla f(x_{k})\|^{2}\le 2(f(x_{0})-f^{*})/(\alpha T)$.

\item \textbf{Constant Noise ($\sigma_{k}^{2}\equiv\sigma^{2}$).}
Then $\sum_{k=0}^{T-1}\sigma_{k}^{2}=T\sigma^{2}$ and (4) becomes
\[
\frac{1}{T}\sum_{k=0}^{T-1}\E\|\nabla f(x_{k})\|^{2}
\le\frac{2(f(x_{0})-f^{*})}{\alpha T}+L\alpha d\sigma^{2},
\]
i.e. the gradient norm plateaus at $O(\alpha)$.

\item \textbf{High‑dimensional limit.} The variance term scales linearly
with $d$ (see $Ld\alpha c$), reflecting the fact that isotropic Gaussian
perturbations affect each coordinate independently.
\end{enumerate}

%%%%%%%%%%%%%%%%%%%%%%%%%%%%%%%%%%%%%%%%%%%%%%%%%%%%%%%%%%%%%%%%%
\section*{Discussion}
The proof follows the classical smooth‑function descent argument
augmented with explicit handling of a *time‑varying* zero‑mean Gaussian
perturbation.  The key insight is that the total injected energy is
finite because $\sum_{k}\sigma_{k}^{2}<\infty$; consequently the noise does
not prevent the algorithm from driving the expected gradient norm to
zero, albeit at a slightly slower $O(\log T/T)$ rate compared with the
noise‑free case.  This rate matches the best known guarantees for
Langevin‑type algorithms with diminishing temperature schedules and
offers a principled trade‑off between exploration (early large noise)
and exploitation (later vanishing noise).

\end{document}